% GENERATED FILE. Edit canonical lessons, then run npm run generate:books.
% canonical-source-hash: fnv1a64:e66f420eb5b6bd9f
% canonical-lessons: FR-C188-pres-de, FR-C188-loin-de, FR-C188-au-fond, FR-C188-a-droite, FR-C188-a-gauche

\chapter{Near to, Far from, At the bottom, On the right, On the left}
\label{ch:fr-near-to-far-from-at-the-bottom-on-the-right-on-the-left}
\hlchaptermodality{188}

\begin{chapteropening}
\textbf{By the end of this chapter:} \emph{I can say near to, far from, at the bottom, on the right and on the left.}

Complete the last lesson of chapter 188: I can say near to, far from, at the bottom, on the right and on the left.
\end{chapteropening}

\section[près de]{près de — near to}

\label{lesson:FR-C188-pres-de}



\begin{quote}
Before the new one: say the French for until, then the French for since.
\end{quote}

\subsection*{You'll want to know: près de}
\textbf{près de} — \textquotedblleft{}near to\textquotedblright{}.

\begin{cousinweb}
\textbf{près de} — \emph{près}, \textquotedblleft{}near\textquotedblright{}, and \emph{de}, \textquotedblleft{}of\textquotedblright{}. \emph{Près} $\leftarrow$ Latin \emph{presse}, \textquotedblleft{}closely\textquotedblright{} (EtymDB 2.1), and \emph{de} $\leftarrow$ Latin \emph{dē}.
\end{cousinweb}

\begin{grammarlens}[title={What you've built}]
One more word for this chapter.
\end{grammarlens}

\subsection*{Your turn}
\begin{itemize}
\raggedright
  \item \emph{Say it:} \emph{près de}
  \item \emph{Say it:} \emph{près de}, once more
  \item \emph{Say it:} say \emph{près de}
  \item \emph{Recall:} say the French for until, then the French for since, then say \emph{près de} again
\end{itemize}

\begin{tcolorbox}[breakable,colback=teal!4,colframe=teal!35!black,title={Before you move on}]
What does près de mean? (\textquotedblleft{}Near to\textquotedblright{}.) Say it once more.
\end{tcolorbox}
\section[loin de]{loin de — far from}

\label{lesson:FR-C188-loin-de}



\begin{quote}
Before the new one: say the French for since, then the French for near to.
\end{quote}

\subsection*{You'll want to know: loin de}
\textbf{loin de} — \textquotedblleft{}far from\textquotedblright{}.

\begin{cousinweb}
\textbf{loin de} — \emph{loin}, \textquotedblleft{}far\textquotedblright{}, and \emph{de}, \textquotedblleft{}from\textquotedblright{}. \emph{Loin} comes from Latin \emph{longē}, \textquotedblleft{}far off\textquotedblright{}, and \emph{de} from \emph{dē}.
\end{cousinweb}

\begin{grammarlens}[title={What you've built}]
One more word for this chapter.
\end{grammarlens}

\subsection*{Your turn}
\begin{itemize}
\raggedright
  \item \emph{Say it:} \emph{loin de}
  \item \emph{Say it:} \emph{loin de}, once more
  \item \emph{Say it:} say \emph{loin de}
  \item \emph{Recall:} say the French for since, then the French for near to, then say \emph{loin de} again
\end{itemize}

\begin{tcolorbox}[breakable,colback=teal!4,colframe=teal!35!black,title={Before you move on}]
What does loin de mean? (\textquotedblleft{}Far from\textquotedblright{}.) Say it once more.
\end{tcolorbox}
\section[au fond]{au fond — at the bottom, at the back}

\label{lesson:FR-C188-au-fond}



\begin{quote}
Before the new one: say the French for near to, then the French for far from.
\end{quote}

\subsection*{You'll want to know: au fond}
\textbf{au fond} — \textquotedblleft{}at the bottom, at the back\textquotedblright{}.

\begin{cousinweb}
\textbf{au fond} — \emph{au}, \textquotedblleft{}at the\textquotedblright{}, and \emph{fond}, \textquotedblleft{}bottom\textquotedblright{}, $\leftarrow$ Latin \emph{fundus}, \textquotedblleft{}bottom\textquotedblright{}. (EtymDB 2.1)
\end{cousinweb}

\begin{grammarlens}[title={What you've built}]
One more word for this chapter.
\end{grammarlens}

\subsection*{Your turn}
\begin{itemize}
\raggedright
  \item \emph{Say it:} \emph{au fond}
  \item \emph{Say it:} \emph{au fond}, once more
  \item \emph{Say it:} say \emph{au fond}
  \item \emph{Recall:} say the French for near to, then the French for far from, then say \emph{au fond} again
\end{itemize}

\begin{tcolorbox}[breakable,colback=teal!4,colframe=teal!35!black,title={Before you move on}]
What does au fond mean? (\textquotedblleft{}At the bottom, at the back\textquotedblright{}.) Say it once more.
\end{tcolorbox}
\section[à droite]{à droite — on the right}

\label{lesson:FR-C188-a-droite}



\begin{quote}
Before the new one: say the French for far from, then the French for at the bottom.
\end{quote}

\subsection*{You'll want to know: à droite}
\textbf{à droite} — \textquotedblleft{}on the right\textquotedblright{}.

\begin{cousinweb}
\textbf{à droite} — \emph{à}, \textquotedblleft{}on\textquotedblright{}, and \emph{droite}, \textquotedblleft{}the right\textquotedblright{}. \emph{Droite} $\leftarrow$ Latin \emph{dīrēctus}, \textquotedblleft{}straight\textquotedblright{} (Lewis \& Short), through Old French \emph{dreit}; \emph{à} $\leftarrow$ Latin \emph{ad}.
\end{cousinweb}

\begin{grammarlens}[title={What you've built}]
One more word for this chapter.
\end{grammarlens}

\subsection*{Your turn}
\begin{itemize}
\raggedright
  \item \emph{Say it:} \emph{à droite}
  \item \emph{Say it:} \emph{à droite}, once more
  \item \emph{Say it:} say \emph{à droite}
  \item \emph{Recall:} say the French for far from, then the French for at the bottom, then say \emph{à droite} again
\end{itemize}

\begin{tcolorbox}[breakable,colback=teal!4,colframe=teal!35!black,title={Before you move on}]
What does à droite mean? (\textquotedblleft{}On the right\textquotedblright{}.) Say it once more.
\end{tcolorbox}
\section[à gauche]{à gauche — on the left}

\label{lesson:FR-C188-a-gauche}



\begin{quote}
Before the new one: say the French for at the bottom, then the French for on the right.
\end{quote}

\subsection*{You'll want to know: à gauche}
\textbf{à gauche} — \textquotedblleft{}on the left\textquotedblright{}.

\begin{cousinweb}
\textbf{à gauche} — \emph{à}, \textquotedblleft{}on\textquotedblright{}, and \emph{gauche}, \textquotedblleft{}the left\textquotedblright{}. \emph{Gauche} comes from Old French \emph{gauchir}, \textquotedblleft{}to warp\textquotedblright{} (EtymDB 2.1), probably Frankish further back; \emph{à} $\leftarrow$ Latin \emph{ad}.
\end{cousinweb}

\begin{grammarlens}[title={What you've built}]
One more word for this chapter.
\end{grammarlens}

\subsection*{Your turn}
\begin{itemize}
\raggedright
  \item \emph{Say it:} \emph{à gauche}
  \item \emph{Say it:} \emph{à gauche}, once more
  \item \emph{Say it:} say \emph{à gauche}
  \item \emph{Recall:} say the French for at the bottom, then the French for on the right, then say \emph{à gauche} again
\end{itemize}

\begin{tcolorbox}[breakable,colback=teal!4,colframe=teal!35!black,title={Before you move on}]
What does à gauche mean? (\textquotedblleft{}On the left\textquotedblright{}.) Say it once more.
\end{tcolorbox}
